\documentclass[11pt,a4paper]{article}

\usepackage[utf8]{inputenc}
\usepackage[T1]{fontenc}
\usepackage{lmodern}
\usepackage[margin=1in]{geometry}
\usepackage{amsmath,amssymb,amsfonts}
\usepackage{booktabs}
\usepackage{array}
\usepackage{hyperref}
\usepackage{url}
\usepackage{microtype}
\usepackage{caption}

\hypersetup{
    colorlinks=true,
    linkcolor=blue,
    citecolor=blue,
    urlcolor=blue,
    pdftitle={Learned Continuous Soft Mixture of Low-Rank Adapters for Multi-Domain Fine-Tuning on Constrained Hardware},
    pdfauthor={Burak Ozturk}
}

\title{\textbf{Learned Continuous Soft Mixture of Low-Rank Adapters for Multi-Domain Fine-Tuning on Constrained Hardware}}

\author{
    \textbf{Burak \"Ozt\"urk} \\
    Independent Researcher, SODALS Project \\
    \texttt{burak.ozturk@sodals-research.org}
}

\date{October 2026 \\ \vspace{0.5em} \small Archive Identifier: \texttt{arxiv/sodals-v2.1-soft-moa-20261014}}

\begin{document}

\maketitle

\begin{abstract}
Fine-tuning foundation language models across diverse mathematical and scientific disciplines on commodity consumer hardware (e.g., NVIDIA RTX 4060, 8GB VRAM) faces a fundamental dilemma: training a single monolithic low-rank adapter (LoRA) or jointly training multiple domain adapters induces the \textbf{Subspace Interference Effect}, where cross-subspace covariance perturbations ($\mathbf{A}_1 \mathbf{A}_2^T \neq 0$) catastrophically erode delicate representations acquired during pretraining. In this work, we propose \textbf{Continuous Soft Mixture-of-Adapters (Soft-MoA)}, a parameter-efficient multi-domain serving framework governed by a lightweight \textbf{Gaussian Phase Router (GPR)}.

GPR extracts intermediate hidden representations from late transformer layers (Layer 32), projecting activations onto calibrated domain anchor vectors ($\tau = 0.004913$, Fisher separation ratio $= 15.30$) to compute a scalar domain divergence $z_i$. This dynamically modulates the weight updates of isolated low-rank adapters across all 128 projection layers via a continuous Sigmoid gate:
\[
\Delta \mathbf{W}(x) = (1 - \alpha) s_1 \mathbf{W}_1 x + \alpha s_2 \mathbf{W}_2 x, \quad \alpha = \sigma(z_i / T)
\]
Evaluated on a comprehensive 170-case multi-domain benchmark comprising 150 MATH-500 competition problems (Domain 1) and 20 canonical Abstract Algebra and Lie Theory invariants (Domain 2), our findings establish:
(1) \textbf{The Subspace Interference Law:} End-to-end joint training of multiple adapters across disjoint domains incurs cross-subspace corruption, erasing delicate frontier representations (the dimension of exceptional Lie algebra $E_7 = 133$ collapsed to $70$, and $\mathfrak{sp}(6) = 21$ degraded to $1$);
(2) \textbf{Inference-Time Blending Superiority:} Maintaining domain adapters in strict physical isolation during training and blending them continuously only at inference time (\textbf{Untrained Soft Blender}) achieves peak overall accuracy (\textbf{138 / 170, 81.18\%}), matching discrete regex rule routing while eliminating brittle manual heuristics; and
(3) \textbf{Hardware Efficiency:} The entire two-adapter ensemble executes concurrently within 7,490 MB VRAM on an 8GB laptop GPU, incurring only 1.3 ms auxiliary routing latency per query.
\end{abstract}

\section{Introduction}

Parameter-efficient fine-tuning via Low-Rank Adaptation (LoRA)~\cite{hu2022lora}, combined with 4-bit NormalFloat quantization (QLoRA)~\cite{dettmers2023qlora}, has become the standard for specializing large language models under strict computational and memory budgets. By decomposing weight updates into low-rank factorization matrices $\Delta \mathbf{W} = (\alpha / r) \mathbf{B}\mathbf{A}$, LoRA drastically curtails trainable parameters and optimizer memory footprints.

However, when language models must concurrently master multiple disparate cognitive competencies---such as high-school contest mathematics (e.g., MATH-500~\cite{hendrycks2021math}) and advanced graduate-level abstract algebra (e.g., Lie theory invariants and root systems)---single low-rank adapters exhibit severe capacity ceilings. While dynamic adapter composition frameworks like LoraHub~\cite{huang2023lorahub} explore linear combinations of adapters, practitioners typically resort to one of two suboptimal strategies:
\begin{enumerate}
    \item \textbf{Joint Multi-Task Fine-Tuning:} Concatenating heterogeneous domain data into a shared gradient trajectory. While superficially attractive, we demonstrate empirically that joint training induces cross-subspace covariance interference ($\mathbf{A}_1 \mathbf{A}_2^T \neq 0$), precipitating catastrophic forgetting of fragile, highly localized representations.
    \item \textbf{Discrete Hard Routing:} Training disjoint domain adapters and switching between them using external text classifiers or regular expression pattern matchers. While discrete routing preserves domain isolation, manual regex parsers are brittle, cannot handle queries exhibiting mixed domain semantics, and introduce orchestration complexity.
\end{enumerate}

To resolve this dilemma, we introduce \textbf{Continuous Soft Mixture-of-Adapters (Soft-MoA)}. Our architecture decouples domain learning into two phases:
\begin{itemize}
    \item \textbf{Offline Isolated Specialization:} Domain cartridges ($\mathbf{W}_1, \mathbf{W}_2$) are trained independently with zero gradient cross-talk.
    \item \textbf{Inference-Time Continuous Soft Blending:} A non-invasive, parameter-efficient Gaussian Phase Router (GPR) probes Layer-32 activations and smoothly interpolates adapter contributions via a temperature-scaled Sigmoid gate.
\end{itemize}

\section{Architecture \& Mathematical Formulation}

\subsection{The Gaussian Phase Router (GPR)}

Earlier experiments that attempted online backpropagation into base weights during routing frequently triggered activation instability. In Soft-MoA, the GPR operates strictly as an \textit{observational readout} at Layer 32 of the frozen Qwen 3.5 9B (NF4) backbone.

Given intermediate hidden activation $\mathbf{h}_{32} \in \mathbb{R}^d$ at the final prompt token, the router projects $\mathbf{h}_{32}$ against calibrated domain centroids $\mathbf{c}_1, \mathbf{c}_2 \in \mathbb{R}^d$:
\begin{equation}
z_i = \frac{\langle \mathbf{h}_{32}, \mathbf{c}_2 - \mathbf{c}_1 \rangle}{\|\mathbf{c}_2 - \mathbf{c}_1\|_2} - \tau
\end{equation}
where $\tau = 0.004913$ is the decision threshold established by empirical Fisher discriminant analysis, yielding a Fisher separation ratio of $15.30$ between canonical MATH-500 and Abstract Algebra prompts.

\subsection{Continuous Soft Blending Equation}

Rather than selecting a discrete adapter $\arg\max_k z_k$, the Soft LoRA Blender continuously modulates the forward pass across all 128 LoRA projection matrices:
\begin{equation}
\mathbf{y} = \mathbf{W}_{\text{base}} \mathbf{x} + (1 - \alpha) s_1 \mathbf{B}_1 \mathbf{A}_1 \mathbf{x} + \alpha s_2 \mathbf{B}_2 \mathbf{A}_2 \mathbf{x}
\end{equation}
where:
\begin{equation}
\alpha = \sigma\left(\frac{z_i}{T}\right) = \frac{1}{1 + \exp(-z_i / T)}, \quad T = 0.03
\end{equation}
At temperature $T = 0.03$, the gating function provides a smooth, differentiable sigmoid transition across domain boundaries while maintaining sharp classification margins for unambiguous queries (operating functionally as a soft-margin switch where $|z_i| > 0.1 \implies \alpha \approx 0$ or $1$). Generation repetition penalty is simultaneously modulated to reflect domain requirements:
\begin{equation}
\text{penalty}(\alpha) = (1 - \alpha) \cdot 1.05 + \alpha \cdot 1.00
\end{equation}

\section{Empirical Multi-Domain Benchmark ($N=170$)}

\subsection{Benchmark Composition}

To stress-test multi-adapter interference, we constructed a unified 170-case multi-domain benchmark:
\begin{itemize}
    \item \textbf{Domain 1 (MATH-500 Head-to-Head Subset, $N=150$):} Canonical contest mathematics problems across 7 disciplines (Algebra, Precalculus, Geometry, Number Theory, Counting \& Probability, Intermediate Algebra, Prealgebra).
    \item \textbf{Domain 2 (Abstract Algebra \& Lie Theory Invariants, $N=20$):} Canonical graduate-level algebra problems testing exact Lie algebra dimensions ($\dim F_4 = 52$, $\dim E_7 = 133$, $\dim E_8 = 248$), root counts, and symplectic invariants ($\dim \mathfrak{sp}(6) = 21$).
\end{itemize}
All evaluations were executed under deterministic greedy decoding ($T=0.0$) directly on an NVIDIA GeForce RTX 4060 Laptop GPU (8GB VRAM).

\subsection{Master Factorial Results}

\begin{table}[htbp]
\centering
\caption{Comprehensive 170-Question Multi-Domain Evaluation across 5 Architectural Variants.}
\label{tab:soft_moa}
\small
\begin{tabular}{@{}cllcccc@{}}
\toprule
\textbf{\#} & \textbf{Architecture / Evaluation Arm} & \textbf{Routing Mechanism} & \textbf{Domain 1 (150Q)} & \textbf{Domain 2 (20Q)} & \textbf{Total (170Q)} & \textbf{Latency} \\
\midrule
1 & Regex Router MoA & Discrete Regex Rules & 122 / 150 (81.33\%) & 16 / 20 (80.00\%) & 138 / 170 (81.18\%) & 0.0 ms \\
2 & \textbf{Untrained Soft Blender} & \textbf{Continuous GPR} & \textbf{122 / 150 (81.33\%)} & \textbf{16 / 20 (80.00\%)} & \textbf{138 / 170 (81.18\%)} & \textbf{1.3 ms} \\
3 & Phase 3 (Scratch Joint LoRA) & Differentiable Gate & 111 / 150 (74.00\%) & \textbf{17 / 20 (85.00\%)} & 128 / 170 (75.29\%) & 2.1 ms \\
4 & Phase 3 (Calibrated Anchor) & Continuous GPR & 118 / 150 (78.67\%) & \textbf{17 / 20 (85.00\%)} & 135 / 170 (79.41\%) & 1.8 ms \\
5 & Phase 13 (Warm-Start Joint) & Continuous GPR & 120 / 150 (80.00\%) & \textbf{17 / 20 (85.00\%)} & 137 / 170 (80.59\%) & 1.4 ms \\
\bottomrule
\end{tabular}
\end{table}

\section{Key Scientific Insights}

\subsection{The Subspace Interference Law}

The primary empirical finding from Table~\ref{tab:soft_moa} is that \textbf{joint multi-adapter training systematically degrades general reasoning capacity}. In Phase 3, joint training caused Domain 1 accuracy to drop from $81.33\%$ to $74.00\%$ ($-11$ problems). Even in Phase 13 (warm-starting from frozen anchors), joint training recovered only 120 / 150 questions (gaining 7 problems but breaking 8 previously solved problems). Standard Gram matrix orthogonality regularizers ($\lambda_{\text{ortho}} \| \mathbf{A}\mathbf{A}^T - \mathbf{I} \|_F^2$) enforce internal Stiefel frame isometry within an adapter, but cannot prevent cross-adapter covariance interference ($\mathbf{A}_1 \mathbf{A}_2^T \neq 0$).

\subsection{Erosion of Delicate Frontier Invariants ($E_7 = 133 \to 70$)}

The cost of joint training is vividly illustrated by Question 12 on exceptional Lie algebra $E_7$ ($\dim E_7 = 133$). In Phase 3, the specialized algebra adapter correctly resolved \texttt{133} (PASS) for the first time in project history. However, 76 gradient steps of joint warm-start fine-tuning in Phase 13 erased this delicate representation, collapsing the prediction to \texttt{70} (FAIL) and corrupting the symplectic invariant $\dim \mathfrak{sp}(6) = 21 \to 1$. Furthermore, on Domain 2, the unadapted frozen base model scored 18 / 20 (90.0\%), showing that gradient updates into late-layer subspaces incur a mild capacity penalty (16 / 20 in Soft-MoA) relative to native pretraining representations.

\subsection{Tuned Lens Probing \& Representation Crystallization}

Using layer-wise affine probes (Tuned Lens~\cite{belrose2023tunedlens}), we tracked where domain representations crystallize across the 32-layer transformer backbone. Probing confirms that abstract algebraic invariants do not separate in early layers (Layers 0--16); domain separation crystallizes almost exclusively in \textbf{Layers 28--32}. Consequently, placing the GPR probe at Layer 32 maximizes signal-to-noise ratio ($F = 15.30$) without disturbing intermediate semantic representations.

\section{Hardware Footprint \& Conclusion}

Evaluated on an NVIDIA GeForce RTX 4060 Laptop GPU (8GB VRAM), the 4-bit NF4 base weights consume 5,500 MB, Domain 1 adapter ($r=4$) consumes 18.2 MB, Domain 2 adapter ($r=8$) consumes 36.4 MB, and runtime activations, KV-cache, LM-head buffers, and CUDA context account for $\sim 1,935\text{ MB}$, totaling \textbf{7,490 MB operational VRAM} ($\sim 698\text{ MB}$ physical headroom, zero WDDM paging over PCIe) with a \textbf{1.3 ms} routing latency.

Continuous Soft-MoA demonstrates that when adapting across disjoint domains on constrained hardware, adapters must remain strictly isolated during training and continuously blended only at inference time. In subsequent downstream serving extensions (\texttt{soft\_lora\_blender.py}), this isolation mechanism has been extended so auxiliary single-purpose identity adapters (e.g., orthogonal \texttt{mlp.down\_proj} identity LoRAs) preserve fixed unit scaling independently of the two-expert $(1-\alpha, \alpha)$ reasoning blend. Future work will extend GPR to hierarchical $K$-adapter routing ($K \ge 8$) and formalize PAC-Bayesian generalization bounds on the mixing coefficient $\alpha$.

\begin{thebibliography}{9}

\bibitem{hu2022lora}
E.~J. Hu, Y.~Shen, P.~Wallis, Z.~Allen-Zhu, Y.~Li, S.~Wang, L.~Wang, and W.~Chen, ``LoRA: Low-Rank Adaptation of Large Language Models,'' in \textit{Proc. ICLR}, 2022.

\bibitem{hendrycks2021math}
D.~Hendrycks, C.~Burns, S.~Kadavath, A.~Arora, S.~Basart, E.~Tang, D.~Song, and J.~Steinhardt, ``Measuring Mathematical Problem Solving With the MATH Dataset,'' in \textit{Proc. NeurIPS}, 2021.

\bibitem{belrose2023tunedlens}
N.~Belrose, Z.~Furman, L.~Smith, D.~Halawi, I.~Ostrovsky, L.~McKinney, S.~Biderman, and J.~Steinhardt, ``Eliciting Latent Predictions from Transformers with the Tuned Lens,'' \textit{arXiv preprint arXiv:2303.08112}, 2023.

\bibitem{dettmers2023qlora}
T.~Dettmers, A.~Pagnoni, A.~Holtzman, and L.~Zettlemoyer, ``QLoRA: Efficient Finetuning of Quantized LLMs,'' in \textit{Proc. NeurIPS}, vol.~36, 2023.

\bibitem{huang2023lorahub}
C.~Huang, Q.~Liu, B.~Y. Lin, T.~Pang, C.~Du, and M.~Lin, ``LoraHub: Efficient Cross-Task Generalization via Dynamic LoRA Composition,'' \textit{arXiv preprint arXiv:2307.13269}, 2023.

\end{thebibliography}

\end{document}
